\section{Predictor architecture}
\label{sec:method-architecture}

The M\textsuperscript{3}N predictor consists of two components:
a neural \emph{backbone} that produces per-cell unary
potentials from the input, and a \emph{Markov Network decision
layer} that combines those unary potentials with a learned
pairwise matrix to score complete labelings. This section
pins down the concrete configurations of both components used
in the experiments of Chapter~\ref{chap:experiments}: the
backbones in
Section~\ref{subsec:method-architecture-backbones}, the edge
sets and pairwise parameterization for the MN layer in
Section~\ref{subsec:method-architecture-mn}, and the
composition that ties the two together in the \texttt{M3N}
class in Section~\ref{subsec:method-architecture-full}.

\subsection{Backbone configurations}
\label{subsec:method-architecture-backbones}

The abstract backbone interface of
Section~\ref{subsec:backbone-interface} is realised through the
\texttt{ConfigBackbone} class: each architecture is declared as a
list of layer specifications in YAML, and the per-node forward
pass shares parameters across all nodes. Two concrete architectures are exercised in this
thesis: a linear-only baseline used on symbolic input data,
and a LeNet-5 convolutional network used on every visual
input.

\paragraph{Linear baseline.}
A single linear layer mapping the per-cell one-hot encoding of
the symbolic observation to $K$ unary scores. Two instances are
used in this thesis: $K = 10$ on symbolic HMC and $K = 9$ on
symbolic Sudoku. The resulting predictor has $K \times K$
trainable backbone parameters plus the $K \times K$ pairwise
matrix --- the simplest possible structured predictor, and a
sanity-check that the LeNet-5 + LP-M3N stack ought to match on
symbolic data and beat on visual data.

\paragraph{LeNet-5.}
The LeNet-5 variant of Section~\ref{sec:nn-features} is the
backbone used on every visual benchmark. The exact YAML layer
sequence is shown in Figure~\ref{fig:method-lenet5-yaml}; the only
difference between the Sudoku ($K = 9$) and HMC ($K = 10$)
variants is the output dimension of the final linear layer.

\begin{figure}[htbp]
\begin{verbatim}
backbone:
  type: config
  layers:
    - {type: conv2d, in_channels: 1,  out_channels: 6,   kernel_size: 5, padding: 2}
    - {type: relu}
    - {type: maxpool2d, kernel_size: 2}
    - {type: conv2d, in_channels: 6,  out_channels: 16,  kernel_size: 5}
    - {type: relu}
    - {type: maxpool2d, kernel_size: 2}
    - {type: conv2d, in_channels: 16, out_channels: 120, kernel_size: 5}
    - {type: relu}
    - {type: flatten, start_dim: 1}
    - {type: linear, in_features: 120, out_features: 84}
    - {type: relu}
    - {type: linear, in_features: 84,  out_features: K}
    - {type: Softmax, dim: 1}
\end{verbatim}
\caption[LeNet-5 backbone YAML]{LeNet-5 backbone, excerpt of
\texttt{configs/architectures/sudoku\_lenet5.yaml}. $K = 9$ on
Sudoku and $K = 10$ on visual HMC; everything else is shared.}
\label{fig:method-lenet5-yaml}
\end{figure}

\subsection{Markov Network decision layer}
\label{subsec:method-architecture-mn}

The decision layer takes the unary potentials produced by the
backbone and combines them with a learned pairwise matrix
$W \in \mathbb{R}^{K \times K}$ to score a complete labeling
$y \in \{1, \ldots, K\}^{|V|}$ via the M\textsuperscript{3}N
score of Section~\ref{sec:nn-features},
\begin{equation}
    f(x, y) \;=\; \sum_{v \in V}
    \bigl[ g_\theta(x_v) \bigr]_{y_v}
    \;+\; \sum_{\{u, v\} \in E} W_{y_u,\, y_v}.
    \label{eq:method-mn-score}
\end{equation}
Two design decisions pin down the layer concretely.

\paragraph{Pairwise parameterization.}
$W$ is a single $K \times K$ matrix shared across every edge
of the graph, so the model carries no explicit notion of edge
type --- on the Sudoku graph, the row, column, and box edges
all consult the same pairwise weights. This is what makes the
all-different constraint of
Section~\ref{sec:markov-networks} \emph{learned soft}: the
network has to discover, from data, that the pairwise potential
should penalise equal-label configurations $W_{a, a}$.
Initialisation is Gaussian with standard deviation
\texttt{pairwise.init\_scale}, set to $0.1$ across all
architecture YAMLs --- small enough that early training is
dominated by the unary term (which is easier to learn first),
large enough that $W$ has a non-zero starting gradient signal
at the first training step.

\paragraph{Edge sets.}
The edge list is not baked into the model. It is produced by a
graph constructor at runtime and passed to the score function
and inference decoders alongside the unary potentials. Two
constructors are used:
\begin{itemize}
    \item \texttt{chain\_edges($T$)} returns the chain
    $0\text{--}1\text{--}\cdots\text{--}(T - 1)$ as a
    $[T - 1, 2]$ tensor. Used by the HMC tasks with
    $T = 30$.
    \item \texttt{sudoku\_edges()} returns the Sudoku
    constraint graph: every pair of cells that shares a row,
    column, or $3 \times 3$ box. The result is an $[810, 2]$
    tensor, matching the $|V| = 81$, $|E| = 810$ counts
    reported in Section~\ref{sec:markov-networks}.
\end{itemize}
Decoupling the edge list from the model means the same
\texttt{M3N} instance is reused across both tasks; the
constructor is selected by the architecture YAML's
\texttt{graph} block.

\subsection{Full predictor}
\label{subsec:method-architecture-full}

The two pieces are composed in the \texttt{M3N} class of
\texttt{mnlearn.models}:
\begin{itemize}
    \item Constructor:
    \texttt{M3N(backbone, num\_classes, pairwise\_init\_scale)}
    stores the backbone module and registers the pairwise
    matrix as a learnable \texttt{nn.Parameter}, so the
    optimizer sees both as trainable parameters in a single
    \texttt{model.parameters()} pass.
    \item \texttt{forward($x$)}: runs the backbone and returns
    the unary potentials of shape $[\mathrm{batch}, |V|, K]$.
    This is the entry point used by inference at evaluation
    time.
    \item \texttt{score($u$, $y$, \mathrm{edges})}: computes
    Eq.~\eqref{eq:method-mn-score} for unary potentials $u$, a
    labeling $y$, and an edge list. Differentiable in $u$
    and in the pairwise matrix; not differentiable in $y$ (an
    integer labeling). Used inside the structured-hinge
    subgradient (Section~\ref{subsec:structured-hinge}) and in
    unit tests.
\end{itemize}

The full predictor is constructed from an
\texttt{ArchitectureCfg} through the \texttt{build\_model}
factory of \texttt{mnlearn.models}: the factory builds the
backbone via \texttt{build\_backbone}, the edge list via
\texttt{build\_graph}, and returns the pair
\texttt{(M3N, edges)} ready to be moved to the training device.